\chapter{Estado del arte}

\section{Economía}

Antes de comenzar con términos más técnicos o informáticos, es importante contextualizar este trabajo, para ello se ha realizado una breve introducción a conceptos básicos y necesarios para este trabajo.

\subsection{Mercados financieros}

Los mercados financieros son el lugar en el que oferentes y demandantes compran y venden activos financieros \cite{raisin_mercados}. Son, por tanto, el lugar donde se encuentran quienes necesitan financiación y quienes disponen de dinero y desean obtener una rentabilidad. La función principal de estos mercados es dirigir el ahorro hacia la inversión productiva, facilitando así el crecimiento económico. A su vez, estos mercados siguen el principio de la oferta y la demanda.

Entre las características más relevantes de los mercados financieros, se encuentran la liquidez, que es la facilidad para comprar o vender un activo, luego está la transparencia, que es el acceso a la información para la toma de decisiones y por último, estos mercados financieros se caracterizan por la volatilidad, que se refiere a los constantes cambios en los precios de los activos, impulsados por eventos geopolíticos, el sentimiento colectivo del mercado o una serie de variables macroeconómicas que veremos más adelante. Tal y como señaló Eugene Fama en su formulación de la Hipótesis de los Mercados Eficientes, "los precios de los activos reflejan plenamente toda la información disponible" \cite{fama1970}, lo que en teoría dificultaría su predicción. Sin embargo, la creciente evidencia de anomalías y la irrupción del Machine Learning han cuestionado esta visión tradicional, abriendo la puerta a enfoques predictivos como el que se plantea en este trabajo. Comprender esta volatilidad y predecir el comportamiento de los activos es el gran desafío del análisis financiero y el objetivo que motiva este trabajo.

\subsection{Fondos de inversión y acciones}
Este punto es clave, pues a través de él se pretende explicar cada uno de los elementos que el modelo va a predecir. A continuación, se definen y contextualizan los dos instrumentos financieros sobre los que recae el análisis predictivo de este trabajo, los fondos de inversión y las acciones.

\begin{itemize}
    \item \textbf{Fondos de inversión}
    Respecto a los fondos de inversión, estos podrían definirse de acuerdo al banco Santander de la siguiente forma: "Un fondo de inversión, de inversión colectiva o inversión común es un instrumento que permite al ahorrador acceder al mercado financiero no de una manera directa, sino a través de inversiones diversificadas" \cite{santander_fondos}. Explicando un poco lo anterior, un fondo de inversión agrupa el dinero de varios ahorradores para invertirlo en una cesta de diferentes activos financieros, siendo todo ello gestionado por un equipo de profesionales que a su vez, se encargan de la toma de decisiones. Algunas de las principales ventajas que ofrecen los fondos de inversión, son la diversificación, al permitir acceder a múltiples mercados y clases de activos con un bajo coste, otra ventaja es la gestión profesional, ya que cuentan con especialistas que buscan maximizar la rentabilidad según el perfil de riesgo y por último, cabe mencionar también la transparencia, ya que estos fondos están supervisados por organismos como la Comisión Nacional del Mercado de Valores (CNMV). En cuanto a los tipos, podemos diferenciar tres: el fondo de renta variable, el de renta fija y el mixto. El presente trabajo se centra en el análisis y predicción del rendimiento de este tipo de instrumentos financieros, empleando para ello técnicas de Machine Learning y series temporales, todo ello con el objetivo de predecir su comportamiento futuro a distintos horizontes temporales.

    \item \textbf{Acciones}
    Al igual que con los fondos de inversión, a través del presente trabajo se pretende calcular la evolución temporal de las acciones de cualquier tipo de compañía. Pero antes de nada, ¿qué son las acciones? ¿En qué se basa su evolución? Una acción es un título que representa una parte del capital social de una empresa. Es por ello que todo aquel que posee acciones de una compañía es, por consiguiente, copropietario de dicha compañía en un porcentaje proporcionalmente relacionado al número de acciones que posea, y es que como bien señala el inversor Warren Buffett, "una acción es un instrumento que te permite poseer una pequeña parte de una gran empresa". En lo que respecta a ser accionista, esto supone un conjunto de derechos como el derecho a recibir parte de los beneficios mediante el reparto de dividendos, el derecho a asistir a las juntas de accionistas y ejercer el voto, el derecho a recibir información cada cierto tiempo sobre la evolución de la compañía y el derecho a recibir la parte proporcional del capital en caso de que la compañía desapareciese. Entrando en el tema de los precios, el valor de una acción no es fijo, sino que varía constantemente en función de la ley de la oferta y la demanda, que depende de varios factores como pueden ser la confianza de los inversores, los resultados de la empresa, las perspectivas de su sector y, de manera muy relevante para este trabajo, de la situación económica general y las variables macroeconómicas.
\end{itemize}

\subsection{Variables macroeconómicas}
Primero de todo, es importante entender que es la macroeconomía, se trata de la rama de la ciencia económica que estudia el funcionamiento de una economía en su conjunto, analizando fenómenos agregados como la producción, el empleo o la inflación \cite{unir_macroeconomia}. En lo que respecta a las variables macroeconómicas, estas son indicadores que miden estos fenómenos, lo que permite entender el contexto en el que se desenvuelven los mercados financieros.

El contexto macroeconómico en el que se va  situar este trabajo es en gran parte estadounidense, debido a que el mercado bursátil de Estados Unidos, representa una gran parte del mercado bursátil global, con valores cercanos al 50\%, ejerciendo una influencia dominante sobre el resto de los mercados internacionales, sin embargo, también se tendrán en cuenta variables europeas. Respecto a las variables seleccionadas para el dataset preanálisis, se encuentran las siguientes.
\begin{table}[H]
\centering
\begin{tabular}{|p{4cm}|p{10cm}|}
\hline
\textbf{Variable} & \textbf{Justificación} \\
\hline
S\&P 500 & Captura el movimiento agregado del mercado de renta variable estadounidense \\
\hline
VIX (Índice de volatilidad) & Refleja la incertidumbre y el riesgo percibido por los inversores \\
\hline
Petróleo & Materia prima clave, cuyo coste impacta directamente en los gastos de empresas e inflación \\
\hline
Oro & Considerado activo refugio, en épocas inflacionarias su precio sube \\
\hline
Tipos de Interés Fed & Principal referencia de política monetaria en EE.UU. \\
\hline
Tipos de Interés BCE & Principal referencia de política monetaria en la zona euro \\
\hline
IPC EE.UU. & Mide el aumento de precios y condiciona la política monetaria en EE.UU. \\
\hline
IPC Europa & Mide el aumento de precios y condiciona la política monetaria en la zona euro \\
\hline
Curva de tipos 10Y2Y & Sirve como predictor de recesiones \\
\hline
Desempleo en EE.UU. & Indica el estado del mercado laboral, influye fuertemente en la economía \\
\hline
Tipo de cambio Euro/Dólar & Divisas principales a nivel global \\
\hline
\end{tabular}
\label{tab:variables_macro}
\caption{Variables macroeconómicas seleccionadas y su justificación}
\end{table}

A parte de las variables mencionadas anteriormente, también se consideró añadir el PIB, sin embargo, pese a ser la variable macroeconómica más representativa del crecimiento económico de un país, el hecho de que se publique trimestralmente y con retraso, hace que no sea una variable útil para modelos predictivos a corto-medio plazo 

\subsection{Contextos económicos}
A su vez, es importante introducir los diferentes contextos económicos que se pueden llegar a encontrar ya que los mercados financieros se pueden llegar a ver muy influenciados por dichos contextos y es que a lo largo de los últimos 22 años, la economía mundial ha atravesado diferentes contextos económicos, cada uno con unas características diferenciales y que tienen un gran impacto, ya sea positivo o negativo en los mercados y las variables macroeconómicas. 

El primer contexto que se diferencia es el de \textbf{expansión económica}, caracterizado por un crecimiento sostenido de la actividad económica, bajada de la tasa de desempleo y VIX bajo y rentabilidades altas. El segundo contexto es totalmente lo opuesto, se trata de \textbf{crisis} que implican una dura bajada de la actividad económica, aumento del desempleo y fuertes caídas en los mercados financieros, un ejemplo cercano sería la crisis de 2008 o la del COVID-19.  El tercer contexto es el de \textbf{alta inflación}, caracterizado por la fuerte subida en los precios, provocando con ello la subida de los tipos de interés. Todo lo contrario a la inflación, es la \textbf{deflación}, caracterizada por la bajada de tipos de interés, impulsando el consumo, la inversión, y los mercados. La identificación del contexto económico vigente en cada periodo resulta fundamental para los modelos predictivos, ya que el comportamiento de los fondos de inversión difiere significativamente entre épocas de expansión, crisis o inflación. Este análisis se retomará en el apartado de clustering del Análisis Exploratorio de Datos (EDA), donde se clasificarán los diferentes periodos temporales tratará más en detalle.

\section{Series temporales}
Este punto es uno de los más importantes, ya que se explica con detalle el concepto de serie temporal. Su relevancia radica en que el dataset final que se proporcionará al modelo será una serie temporal multivariante de frecuencia mensual. Por ello, resulta necesario explicar todos los conceptos relacionados, así como el análisis que se pretende seguir para la selección de variables.

\subsection{Conceptos de series temporales}
Una serie temporal se define como una colección de observaciones de una variable recogidas a lo largo del tiempo, normalmente en instantes equiespaciados. En estas series, las observaciones sucesivas no son independientes entre sí; existe una dependencia temporal intrínseca donde un valor está relacionado con su pasado. Como señalaron George Box y Gwilym Jenkins: ``en una serie temporal, el orden de las observaciones es esencial, ya que los valores pasados contienen información sobre el futuro'' \cite{box_jenkins1970}. Un ejemplo claro son los tipos de interés de la FED: un valor en un mes determinado está condicionado por el del mes anterior, ya que estas variables no suelen presentar cambios bruscos de forma aislada. Por ello, el análisis debe realizarse respetando estrictamente el orden cronológico.

En cuanto a su clasificación, las series pueden ser deterministas o estocásticas. Las primeras son aquellas cuyos valores futuros pueden predecirse de forma exacta mediante una ley matemática. Por el contrario, en las estocásticas los valores futuros solo pueden conocerse a partir de una función de probabilidad. En el ámbito financiero, la mayoría de las series presentan un comportamiento principalmente estocástico, lo que justifica el uso de modelos de \textit{Machine Learning} para su análisis, al no existir una fórmula determinista que permita conocer el valor futuro con certeza.

Continuando con su estructura, toda serie temporal puede descomponerse normalmente en cuatro componentes:

\begin{description}
    \item[Tendencia:] Cambio suave y a largo plazo en el nivel medio de la serie, que representa su dirección general (ascendente, descendente o lateral).
    \item[Estacionalidad:] Variaciones periódicas y predecibles que se repiten en intervalos de tiempo fijos (anuales, trimestrales, etc.).
    \item[Componente cíclico:] Variaciones de mayor duración que la estacionalidad, pero sin un periodo fijo, como los ciclos económicos de recesión.
    \item[Componente aleatoria:] Variaciones impredecibles que permanecen una vez eliminados los efectos anteriores.
\end{description}

Es importante introducir aquí el concepto de \textbf{estacionariedad}, que no debe confundirse con la estacionalidad. Una serie es estacionaria cuando su media y su variabilidad permanecen constantes a lo largo del tiempo. Por el contrario, las series no estacionarias presentan propiedades estadísticas que cambian cronológicamente, lo que suele requerir transformaciones previas para su modelado.

En la Figura \ref{fig:ibex35} se muestra un ejemplo de una serie no estacionaria y no cíclica, con una tendencia ascendente. La imagen detalla la evolución del índice IBEX35 en el periodo comprendido entre 2021 y 2026.

\begin{figure}[H]
    \centering
    \includegraphics[width=0.8\textwidth]{imagenes/Evolucion_IBEX35.png}
    \caption{Evolución del índice IBEX35 (2021-2026). Fuente: BME Exchange}
    \label{fig:ibex35}
\end{figure}

Finalmente, alineando esta explicación a este trabajo, el dataset final que se proporcionará a los modelos predictivos es una serie temporal multivariante de frecuencia mensual, que abarca desde diciembre de 2003 hasta septiembre de 2025. El hecho de que sea multivariante implica que, para cada mes, se dispone de múltiples variables como, por ejemplo, todas las variables macroeconómicas mencionadas anteriormente. Este enfoque multivariante es especialmente relevante en el dominio financiero, donde las correlaciones entre diferentes variables son complejas. Por su parte, la frecuencia mensual ofrece un equilibrio adecuado entre la disponibilidad de datos históricos y la capacidad predictiva a horizontes de corto, medio y largo plazo, que son los tres escenarios contemplados en los objetivos de este trabajo.

\subsection{Técnicas de análisis de series temporales}
\begin{itemize}
    \item \textbf{Análisis de series temporales}
    Una vez definidos los conceptos fundamentales, es necesario abordar las técnicas de análisis de las series temporales, el cual consta de diferentes etapas y cada una de las cuales aporta información importante sobre el comportamiento de los datos. Para este trabajo, el análisis que se va a llevar a cabo es un EDA, que incluye la recopilación, el preprocesamiento y la visualización de las series, pudiendo así entender mejor cada una de las variables. Tras el EDA, iría la selección del modelo, entrenamiento y evaluación, pero eso se verá más tarde.

    Antes de detallar las herramientas concretas, conviene distinguir entre dos enfoques posibles en el análisis de series temporales. Por un lado, el análisis univariante se centra en el estudio de una única variable a lo largo del tiempo, examinando sus propias características, su evolución, su memoria temporal y su capacidad de predicción a partir de sus propios valores pasados. Por otro lado, el análisis multivariante incorpora simultáneamente múltiples variables, lo que permite estudiar las interrelaciones entre ellas y aprovechar la información de unas variables para predecir el comportamiento de otras. Dado que el presente trabajo dispone de un dataset con once variables macroeconómicas además del activo financiero objetivo, el análisis que se llevará a cabo será fundamentalmente multivariante, si bien se complementará con análisis univariantes de cada variable para comprender su comportamiento individual.

    En cuanto a las herramientas más usadas en el análisis de series temporales, una de las principales es el análisis visual, esta herramienta permite comprender una serie temporal de forma mucho más clara que solo viendo datos. Algunos ejemplos de este análisis visual, son los gráficos, que permiten identificar múltiples características de los datos, como la frecuencia de las observaciones, la tendencia creciente o decreciente, la existencia de valores atípicos, la dispersión de los datos alrededor de su media, los cambios de pendiente, la existencia de patrones que se repiten a lo largo del tiempo, los diferentes cuartiles, la media y mediana o características como la curtosis y skewness, ambas son muy útiles a la hora de entender comportamientos de variables. Para acabar, algunos ejemplos de gráficos, podrían ser gráficos de barras, boxplots o de dispersión entre otros

    Otra de las herramientas a destacar, es la función de autocorrelación, la cual es muy útil en la exploración de series temporales ya que esta función analiza los datos de series temporales para buscar correlaciones en los valores de distintos puntos de una serie temporal, es decir, mide la correlación de un valor consigo mismo \cite{ibm_autocorrelacion}. La representación gráfica de la función de autocorrelación, se denomina correlograma y  resulta útil para identificar estacionalidad, tendencias y otros patrones en los datos, un claro ejemplo de correlograma es la Figura \ref{fig:Correlograma}. La herramienta de autocorrelación, es fundamental para determinar la memoria temporal de cada variable y para ver si los valores pasados contienen información relevante acerca del valor futuro.

    \begin{figure}[H]
    \centering
    \includegraphics[width=0.8\textwidth]{imagenes/CorrelogramaEjemplo.png}
    \caption{Ejemplo de un correlograma de una variable. Fuente: RPubs}
    \label{fig:Correlograma}
    \end{figure}
    Al igual que la autocorrelación, también existe la correlación de Pearson, la cual es muy útil en el análisis de series temporales multivariantes ya que esta función analiza la relación lineal entre dos variables diferentes, es decir, mide cómo se comporta una variable cuando la otra aumenta o disminuye. La representación gráfica de la correlación entre múltiples variables se denomina matriz de correlación y resulta útil para identificar relaciones directas (cuando una variable sube, la otra también sube), relaciones inversas (cuando una variable sube, la otra baja) o ausencia de relación lineal entre las variables. La herramienta de correlación es fundamental para entender cómo se relacionan las distintas variables macroeconómicas entre sí y, especialmente, cómo se relacionan con el activo financiero que se pretende predecir, lo que permite identificar qué variables pueden tener mayor poder explicativo sobre el comportamiento del target.

    Por último, otra herramienta de gran utilidad es el clustering, que permite agrupar las observaciones en función de su similitud. Aplicado a series temporales, esta técnica permite segmentar el período de estudio en diferentes fases o regímenes económicos, identificando automáticamente momentos en los que el comportamiento de las variables presenta patrones similares, como períodos de expansión, crisis o recuperaciones.

    En el contexto de este trabajo, y dado que el dataset final es una serie temporal multivariante de frecuencia mensual que abarca desde diciembre de 2003 hasta septiembre de 2025, el análisis de series temporales se realizará combinando todas las técnicas mencionadas. En primer lugar, se realizará un análisis descriptivo de cada variable macroeconómica por separado, examinando su evolución temporal a lo largo del tiempo, su distribución estadística, sus valores extremos y su volatilidad en diferentes subperiodos. Adicionalmente, se estudiará la autocorrelación de cada variable mediante las funciones ACF, lo que permitirá determinar su memoria temporal y, por tanto, si resulta necesario incluir sus valores pasados como predictores en los modelos. Finalmente, dado el carácter multivariante del dataset, el análisis no se limitará al estudio individual de cada variable, sino que se extenderá a un análisis multivariante que explore las relaciones entre las distintas variables macroeconómicas y su interacción con el activo financiero objetivo. En este sentido, se calcularán las correlaciones entre todas las variables para identificar posibles relaciones lineales y ver si algunas son redundantes entre sí, a su vez,  se aplicará un análisis de componentes principales (PCA) para reducir la dimensionalidad del sistema y comprender cuántas dimensiones explican la mayor parte de la varianza del conjunto de datos, y por último, se llevará a cabo un análisis de clustering con el objetivo de poder diferenciar los diferentes ciclos a lo largos del tiempo de estudio. 

    Todos estos análisis, que se presentan en detalle en el capítulo de Análisis Exploratorio de Datos (EDA), constituyen la base sobre la que se construirán los modelos predictivos en la fase posterior del trabajo.

    \item \textbf{Selección de variables}
    Una vez comprendida la naturaleza de las series temporales y las herramientas para su análisis, el siguiente paso clave es la selección de las variables más relevantes y que conformarán el dataset definitivo que le será entregado a los modelos predictivos. Como se ha visto anteriormente, en este trabajo se dispone de once variables macroeconómicas además de las variables propias del activo financiero. Sin embargo, no todas estas variables aportan la misma cantidad de información a la hora de realizar la predicción de un activo o fondo de inversión, ya que algunas pueden ser redundantes entre sí, otras pueden contener ruido y otras pueden tener una relación no lineal el target. Por esto mismo, la selección de variables es importante, ya que permite por un lado mejorar el potencial predictivo del dataset eliminando variables redundantes y que introducen ruido, y por otro lado, la selección de variables permite reducir la dimensión del dataset, consiguiendo así un menor coste computacional y una mayor interpretabilidad de los resultados. Es importante destacar que todo este proceso de selección se realiza antes que el entrenamiento de los modelos finales, asegurando que el dataset que se entregue a los modelos predictivos, sea ya el conjunto definitivo y optimizado.

    Entrando ahora a los diferentes métodos de selección de variables, existen diferentes formas de métodos de selección de variables, entre los que destacan los métodos de filtrado, los métodos envolventes y los métodos embebidos. Los métodos de filtrado evalúan la relevancia de cada variable de forma independiente a cualquier modelo predictivo posterior, basándose únicamente en propiedades estadísticas de los datos. Estos métodos son computacionalmente eficientes y escalables, aunque no consideran las interacciones entre variables. Los métodos envolventes, por el contrario, evalúan subconjuntos de variables entrenando un modelo predictivo para cada subconjunto, lo que los hace más precisos pero también mucho más costosos computacionalmente. Por último, los métodos embebidos realizan la selección de variables durante el propio proceso de entrenamiento de un modelo. En este trabajo, se emplearán principalmente métodos de filtrado, junto con un análisis de correlación para detectar redundancias.

    Entre los métodos de filtrado, destacan la entropía de datos y la información mutua. La entropía mide el grado de desorden asociado a una variable aleatoria. A la hora de seleccionar variables, si una variable presenta una entropía elevada, significa que dicha variable, presenta una gran variedad de valores y por consiguiente, contiene más información relevante. Sin embargo, si la entropía es baja, entonces significa que no aporta mucha información al dataset y entonces se podría eliminar del dataset. En el caso de este trabajo, se aplicará el estudio de la entropía a cada una de las variables del dataset y se seleccionarán aquellas que realmente aporten información, es decir, que tengan una entropía alta.

    Por otro lado, la información mutua, mide la cantidad de información que una variable contiene sobre otra. Es importante no confundirlo con lo correlación de Pearson mencionada en las técnicas de análisis, ya que dicha correlación, sólo captura relaciones lineales, mientras que la información mutua es capaz de detectar cualquier tipo de dependencia estadística, ya sea lineal o no lineal. Dicha propiedad es muy útil a la hora de hacer este trabajo, ya que las relaciones entre variables macroeconómicas y del activo, no suelen ser lineales. Finalmente, en este trabajo, se calculará la información mutua entre cada variable predictora y la variable objetivo, lo que permitirá ordenar las variables por su capacidad predictiva individual y seleccionar para el dataset definitivo aquellas con mayor información mutua.

    Además de estos métodos de filtrado, se empleará también un análisis de correlación entre las variables del dataset, con el objetivo de detectar posibles redundancias. En caso de que las correlaciones entre variables sean muy altas, entonces se consideraría la eliminación de una de ellas, con el objetivo de no introducir ruido. 

    En lo que respecta a este trabajo, el enfoque práctico que se seguirá para seleccionar las variables finales del dataset seguirá el siguiente proceso. En primer lugar, se analizarán las correlaciones entre las variables predictoras para detectar posibles redundancias. En segundo lugar, se aplicará un filtrado basado en entropía para eliminar aquellas variables con muy baja variabilidad y que no aportan información. En tercer lugar, se calculará la información mutua de cada variable restante con respecto al rendimiento del activo objetivo. Finalmente, con toda esta información, se montará el dataset final, el cual contendrá las variables más relevantes de acuerdo a los filtros anteriores. Este dataset final será el que se entregue a los modelos de Random Forest, Gradient Boosting y LSTM para su entrenamiento y posteriores experimentos.
\end{itemize}

\section{Machine Learning}
Entrando ahora en la parte fundamental del trabajo, los modelos predictivos, los cuales serán los encargados de procesar la información del dataset y generar predicciones. Primero de todo el Machine Learning, el cual es el subconjunto de la Inteligencia Artificial centrado en algoritmos que pueden aprender los patrones de los datos de entrenamiento y, posteriormente, hacer inferencias sobre nuevos datos. Esta capacidad de reconocimiento de patrones permite a los modelos de machine learning tomar decisiones o realizar predicciones sin instrucciones explícitas y codificadas de forma rígida \cite{ibm_machine_learning}. En resumen, un modelo de Machine Learning aprende a partir de ejemplos, identificando relaciones y estructuras que se encuentran en los datos. En un modelo de Machine Learning, cuantos más datos disponga sobre el tema, más precisas serán las predicciones que haga, ya que dispone de más información 

Dentro del Machine Learning, se encuentran varias formas de aprendizaje, siendo el \textbf{aprendizaje supervisado} uno de los más relevantes para este trabajo. En el aprendizaje supervisado, el algoritmo se entrena a partir de un conjunto de datos que contiene tanto las variables de entrada como la variable objetivo que se quiere predecir. En esta forma de aprendizaje, el modelo aprende los patrones a partir de los cuales las entradas producen las salidas, y dichos patrones, los usa luego para predecir la variable objetivo a partir de nuevos datos de entrada, esto se puede ver en la Figura \ref{fig:aprendizaje_super}. Puesto que este trabajo es un problema de regresión donde se pretende predecir una variable continua, entonces el aprendizaje  supervisado es la forma de aprendizaje adecuada. Por último, el aprendizaje supervisado se divide en problemas de clasificación, cuando la variable objetivo es categórica, y problemas de regresión, cuando es continua, siendo este último el caso de este trabajo.

\begin{figure}[H]
\centering
\includegraphics[width=0.7\textwidth]{imagenes/EjemploAprendizajeSuper.png}
\caption{Esquema del funcionamiento del aprendizaje supervisado.}
\label{fig:aprendizaje_super}
\end{figure}

Continuando con el proceso de entrenamiento de un modelo de Machine Learning, este se basa en la optimización de una \textbf{función de pérdida}, la cúal calcula la diferencia entre el valor de la predicción realizada y el valor real. En problemas de regresión, una de las funciones de pérdida más comunes es el error cuadrático medio (MSE), que calcula la media de los cuadrados de las diferencias entre predicciones y valores reales, otros ejemplos de funciones de pérdida, serían el error absoluto medio (MAE) o el error cuadratico medio raiz (RMSE). En Machine Learning, durante el entrenamiento, el objetivo es encontrar los parámetros del modelo que minimicen esta función de pérdida y para lograrlo, se hace uso de un algoritmo denominado \textbf{descenso de gradiente}, el cual calculando el gradiente de la función de pérdida con respecto a los parámetros del modelo, lo que indica la dirección en la que los parámetros deben ajustarse para reducir el error. En cada iteración, los parámetros se actualizan alejándose cada vez más del gradiente, este proceso se repite durante varias iteraciones hasta que el modelo alcanza un error muy bajo.

Otro concepto fundamental en la construcción de modelos predictivos es la reducción de la dimensionalidad. Cuando se trabaja con un conjunto de datos que contiene muchas variables, no todas ellas aportan información relevante para la predicción, algunas pueden ser redundantes entre sí, otras pueden contener principalmente ruido y algunas pueden tener una relación muy débil con el target. Es por ello, que incluir todas estas variables en el modelo, puede provocar sobreajuste o introducir ruido. Por esto mismo, es necesario aplicar técnicas de selección de variables y reducción de dimensionalidad sobre el dataset antes de entrenar los modelos. En este trabajo, esta tarea se ha abordado en profundidad durante la fase de Análisis Exploratorio de Datos (EDA) y se ha explicado en profundidad en el punto anterior, es por ello que aquí no se va a entrar en detalles.

Finalmente, en este trabajo, se trata de resolver un problema multivariante-multistep. Por un lado, es multivariante ya que se como se ha explicado anteriormente, se disponen varias variables para cada instante de tiempo y que actúan a la vez para predecir la variable objetivo. Esta característica exige que los modelos sean capaces de manejar entradas de alta dimensionalidad y de capturar las interacciones y dependencias entre las distintas variables. Por otro lado, es un problema multistep, ya que no se pretende predecir únicamente el valor del mes siguiente, sino que se busca generar predicciones a diferentes horizontes temporales, lo cual añade un punto de complejidad, ya que los errores tienden a acumularse, es por ello que el modelo debe de ser capaz de mantener la precisión.

Los modelos de Machine Learning que se van a usar, son Random Forest y Gradient Boosting. Ambos son modelos basados en árboles de decisión que pertenecen a la familia de \textbf{Ensemble Methods}. Estos métodos combinan múltiples modelos base para obtener un predictor final más robusto y preciso que cualquiera de los modelos individuales. La agregación de varias predicciones reduce la varianza y el error general del modelo, mejorando así su capacidad de generalización. Existen dos enfoques principales, el \textbf{bagging}, donde los modelos se entrenan de forma independiente y en paralelo, un ejemplo será Random Forest, y luego está \textbf{boosting}, en el que los modelos se entrenan de forma secuencial, corrigiendo los errores de los anteriores, un ejemplo de boosting, sería Gradient Boosting.

\subsection{Random Forest}
Random Forest \cite{scikit_rf}, es un algoritmo de aprendizaje automático que combina el resultado de múltiples árboles de decisión para llegar a un resultado único. Puede ser usado tanto en problemas de clasificación como de regresión debido a su flexibilidad. Random Forest es un modelo de árbol de decisión, eso significa que sigue un modelo muy similar al juego de quien el quien, donde a partir de preguntas, se van creando subconjuntos hasta llegar a la respuesta correcta, una definición más técnica sería que árbol de decisión divide recursivamente los datos en subconjuntos más homogéneos mediante nodos de decisión.

Como se ha comentado anteriormente, Random Forest, pertenece a la familia de Ensemble Methods. A su vez, es una extensión del método de bagging, que consiste en seleccionar muestras aleatorias del conjunto de entrenamiento con reemplazo, entrenar un árbol de decisión sobre cada subconjunto y agregar las predicciones promediando en el caso de regresión. Finalmente, este algoritmo de Machine Learning, añade una capa adicional llamada la aleatoriedad de características, que consiste en seleccionar un subconjunto aleatorio de características en cada división del árbol, lo que garantiza que los árboles del bosque no estén correlacionados entre sí, reduciendo la varianza y el riesgo de sobreajuste.

Entre las principales ventajas del Random Forest para este trabajo destacan las siguientes:

\begin{itemize}
    \item \textbf{Menor riesgo de sobreajuste}: al combinar múltiples árboles entrenados sobre diferentes subconjuntos de datos y características, el promediado de sus predicciones reduce la varianza del modelo, evitando que se ajuste excesivamente al ruido de los datos de entrenamiento.
    
    \item \textbf{Flexibilidad para manejar regresiones sin preprocesamiento extensivo}: Random Forest no requiere normalización ni estandarización de las variables, y es robusto frente a valores atípicos y datos faltantes, lo que simplifica el preprocesamiento del dataset.
    
    \item \textbf{Capacidad para evaluar la importancia de las características}: el algoritmo puede medir la contribución de cada variable a la reducción de la impureza en los nodos de los árboles, permitiendo identificar qué variables macroeconómicas son más relevantes para la predicción.
\end{itemize}

Con respecto a las desventajas, este modelo presenta un entrenamiento más lento que otros modelos debido a que tiene que construir múltiples árboles de forma independiente, lo que puede ser costoso computacionalmente si se utiliza un número elevado de árboles o conjuntos de datos muy grandes. Además, su mayor complejidad interpretativa hace que sea más difícil explicar las predicciones individuales que con un único árbol de decisión, aunque las medidas de importancia de variables ayudan a solucionar este problema.

En cuanto a este trabajo, se implementará un \textit{RandomForestRegressor} con scikit-learn \cite{scikit_learn}, ajustando sus hiperparámetros principales como el número de árboles, la profundidad máxima y el número de características por división.
\subsection{Gradient Boosting}

Por otro lado, Gradient Boosting \cite{scikit_gb} es otro Ensemble Method basado en árboles de decisión, pero a diferencia del Random Forest que entrena los árboles de forma independiente y en paralelo, Gradient Boosting los construye de forma secuencial y acumulativa. Pertenece a la familia de métodos de boosting, cuyo objetivo es combinar múltiples modelos débiles para crear un modelo fuerte, donde cada nuevo modelo se centra en corregir los errores cometidos por los anteriores.

Su funcionamiento es el siguiente. Primero, se inicializa el modelo con una predicción base (generalmente la media de la variable objetivo). En cada iteración, se calculan los errores (residuos) del modelo actual y se entrena un nuevo árbol de decisión para predecir esos errores. Este nuevo árbol se añade al modelo multiplicado por un factor \textit{learning rate}, que controla la contribución de cada árbol y evita el sobreajuste. El proceso se repite durante un número determinado de iteraciones o hasta que el error sea suficientemente pequeño.

En cuanto a los hiperparámetros, el \textit{learning rate} es el más importante, ya que un valor bajo hace que el modelo aprenda más lentamente pero de forma más generalizable, debido a que se requieren más árboles para aproximar la función objetivo. Otros hiperparámetros relevantes son la profundidad máxima de cada árbol y la fracción de muestras utilizada en cada iteración.

Entre las principales ventajas del Gradient Boosting para este trabajo destacan las siguientes:

\begin{itemize}
    \item \textbf{Rendimiento predictivo superior}: al construir los árboles de forma secuencial corrigiendo los errores de los anteriores, el modelo puede capturar patrones más complejos y sutiles que otros algoritmos, lo que suele traducirse en una mayor precisión predictiva.
    
    \item \textbf{Capacidad para evaluar la importancia de las características}: al igual que Random Forest, Gradient Boosting puede medir la contribución de cada variable a la reducción del error durante el proceso de entrenamiento, lo que permite identificar qué variables macroeconómicas son más relevantes para la predicción.
\end{itemize}

Con respecto a las desventajas, su entrenamiento es más lento que el de Random Forest debido a su naturaleza secuencial. Además, es más sensible a los hiperparámetros, especialmente al \textit{learning rate} y al número de árboles. Un ajuste inadecuado puede provocar sobreajuste si el aprendizaje es demasiado rápido o infraprendizaje si es demasiado lento. Por ello, el ajuste de hiperparámetros debe realizarse con especial cuidado, utilizando técnicas como validación cruzada para encontrar la configuración óptima.

Finalmente, en este trabajo se implementará un \textit{GradientBoostingRegressor} con scikit-learn \cite{scikit_gb}, ajustando sus hiperparámetros principales mediante validación cruzada.
\section{Deep Learning}
El Deep Learning es un subconjunto del Machine Learning impulsado por redes neuronales cuyo diseño se inspira en la estructura del cerebro humano y que implica el entrenamiento de modelos de al menos 4 capas. A diferencia de los modelos de Machine Learning, que siguen una lógica matemática, las redes neuronales de los modelos de Deep Learning comprenden muchas capas interconectadas de neuronas, que realizan cada una una operación matemática. En estos modelos de Deep Learning, se puede ajustar los pesos y sesgos variables del modelo para así optimizar la red y obtener salidas más precisas.

Continuando ahora con las redes neuronales, estas son un modelo de aprendizaje automático que apila neuronas simples en capas y aprende a reconocer patrones a partir de datos para así asignar entradas a salidas. Los componentes básicos de una red neuronal son una capa de entrada, que contiene las características sin procesar; una o varias capas ocultas, que transforman las entradas en nuevas representaciones mediante funciones de activación no lineales; y una capa de salida, que genera la predicción final. En la Figura \ref{fig:red_neuronal} se pueden ver dichas capas. Se ha mencionado que en las capas ocultas las entradas se transforman con funciones de activación. Algunos ejemplos de estas funciones son la tangente hiperbólica, la unidad lineal rectificada (ReLU) o la función sigmoide. Todas ellas introducen no linealidad en el modelo, permitiendo a la red capturar patrones complejos y relaciones no lineales en los datos.

\begin{figure}[H]
\centering
\includegraphics[width=0.8\textwidth]{imagenes/redNeuronalMemoria.png}
\caption{Estructura de una red neuronal con capas de entrada, ocultas y salida. Fuente: OpenWebinars}
\label{fig:red_neuronal}
\end{figure}

El entrenamiento de una red neuronal se basa en dos algoritmos fundamentales: la retropropagación \cite{ibm_retropropagacion} y el descenso de gradiente \cite{ibm_descenso_gradiente}, que se explicó anteriormente en el apartado de Machine Learning. Profundizando ahora en la retropropagación, esta es un método muy útil para calcular cómo afectan los cambios en cualquier peso o sesgo individual de una red neuronal a la precisión de las predicciones del modelo. Para ello, durante el entrenamiento, se compara la predicción de la red con el valor real mediante una función de pérdida y, a continuación, la retropropagación propaga este error hacia atrás a través de la red desde la capa de salida hasta la capa de entrada, calculando mediante la regla de la cadena del cálculo el gradiente de la función de pérdida con respecto a cada parámetro del modelo. Este gradiente indica en qué dirección deben ajustarse los pesos y sesgos para reducir el error. Posteriormente, un algoritmo de descenso de gradiente utiliza este gradiente para actualizar los parámetros del modelo. Finalmente, este proceso se repite en bucle hasta que el modelo encuentra un conjunto de parámetros que producen predicciones suficientemente precisas.

Sin embargo, es importante señalar que el descenso de gradiente no garantiza la convergencia al mínimo global de la función de pérdida. En la práctica, la función de pérdida de una red neuronal es no convexa y presenta una gran cantidad de mínimos locales y puntos de silla. Es por ello que el algoritmo puede quedar atrapado en un mínimo local, que no necesariamente es la solución óptima, o por el contrario, puede oscilar sin converger si la tasa de aprendizaje no es la adecuada. No existe una única solución, ya que diferentes inicializaciones aleatorias de los pesos pueden llevar a distintos conjuntos de parámetros finales que producen un rendimiento muy parecido. Además, la robustez del proceso depende fuertemente de la elección de la tasa de aprendizaje y de la calidad de los datos. Una tasa de aprendizaje demasiado alta puede provocar divergencia u oscilaciones, mientras que una demasiado baja puede ralentizar excesivamente la convergencia.

Antes de concluir con el Deep Learning es importante mencionar el equilibrio entre sesgo y varianza, que busca reducir el error. El sesgo mide lo lejos que están las predicciones de un modelo de los valores reales, mientras que la varianza captura cuánto varían las predicciones en función de los diferentes datos de entrenamiento. Un modelo con alto sesgo es propenso a un ajuste insuficiente, es decir, no logra capturar los patrones importantes de los datos y da lugar a un modelo demasiado simple. Por el contrario, un modelo con alta varianza es propenso al sobreajuste, lo que significa que se ajusta excesivamente a los datos de entrenamiento, aprendiendo el ruido como si fuera señal y perdiendo capacidad de generalización a nuevos datos. Con el entrenamiento, se busca encontrar el punto óptimo entre ambos errores, minimizando el error de predicción. Algunas técnicas para mantener este equilibrio son la regularización, que añade una penalización a la función de pérdida para desalentar modelos demasiado complejos, también está el \emph{dropout}, que apaga aleatoriamente neuronas durante el entrenamiento para reducir el sobreajuste y por último, la detención temprana, que detiene el entrenamiento cuando la pérdida en validación deja de mejorar.

\subsection{LSTM}
Las redes neuronales recurrentes (RNN) son un tipo de red neuronal diseñada específicamente para trabajar con datos secuenciales, como series temporales. A diferencia de las redes neuronales convencionales, las RNN cuentan con bucles de retroalimentación que permiten que la información se mantenga a lo largo de los pasos temporales, creando una memoria de entradas pasadas. Sin embargo, las RNN convencionales presentan una limitación importante, el problema de la desaparición del gradiente, el cual impide que la red aprenda dependencias a largo plazo.

Para superar este problema, surge la arquitectura Long Short-Term Memory (LSTM), diseñada para capturar dependencias temporales a largo plazo de manera efectiva y es que tal y como mencionan Sepp Hochreiter y Jürgen Schmidhuber "las redes LSTM pueden aprender a puentear intervalos de tiempo de más de 1000 pasos sin perder información relevante" \cite{hochreiter1997}. Esto lo consigue gracias a su estructura interna, que cuenta con tres puertas que regulan el flujo de información, una de las puertas es la \textbf{puerta de olvido}, que decide qué información del pasado debe descartarse, luego está la \textbf{puerta de entrada}, que determina qué nueva información debe almacenarse, y por último, la \textbf{puerta de salida} controla qué información se utiliza para generar la predicción. Gracias a estas puertas, LSTM, es capaz de sobreponerse a la limitación planteada anteriormente y además, permite a la red olvidar información que no es relevante.

En el contexto de este trabajo, la LSTM presenta ventajas bastante interesantes. Una de las ventajas, como se ha mencionado antes, es la gran capacidad de LSTM para capturar dependencias temporales prolongadas. En segundo lugar, LSTM es capaz de manejar el carácter multivariante del problema, pudiendo procesar a la vez todas las variables macroeconómicas y del activo como una secuencia de entrada. Por último, LSTM es capaz de abordar el problema multi-step, generando predicciones a diferentes horizontes temporales y manteniendo la coherencia temporal de las predicciones. En cuanto a las desventajas, LSTM requiere una mayor cantidad de datos para entrenarse adecuadamente, es más costosa computacionalmente y la interpretabilidad de las predicciones es más costosa que en modelos de Machine Learning. 

En este trabajo se implementará una red LSTM utilizando TensorFlow \cite{tensorflow}, configurándola para el problema multivariante-multistep, y se comparará su rendimiento con el de Random Forest y Gradient Boosting, evaluando si su capacidad de memoria a largo plazo ofrece ventajas significativas sobre los métodos basados en árboles para la predicción de rendimientos financieros.

\section{Modelos explicativos}
En este trabajo no solo se pretende realizar predicciones, sino que también se busca entender el porqué de las salidas del modelo, es por ello que se emplearán modelos explicativos para así poder comprender la influencia de cada una de las variables en las predicciones y entender mejor el comportamiento del modelo. Para conseguir esta transparencia, se emplearán dos técnicas de Explainable Artificial Intelligence (XAI) \cite{xai_ibm}, LIME y SHAP, las cuales permitirán interpretar las predicciones del modelo y analizar la influencia de cada variable macroeconómica en los resultados obtenidos.

\subsection{LIME}
LIME \cite{xai_lime} es una técnica explicativa que genera explicaciones locales para predicciones individuales. Su funcionamiento consiste en aproximar el comportamiento del modelo complejo en la vecindad de una instancia específica mediante un modelo más simple e interpretable, como una regresión lineal. De esta forma, para una predicción concreta del rendimiento de un fondo en un mes determinado, LIME permite identificar qué variables macroeconómicas han influido más en esa decisión concreta, ofreciendo una explicación comprensible y localizada. Su principal ventaja es que puede aplicarse tanto a Random Forest, como a Gradient Boosting o a LSTM sin necesidad de modificaciones.

\subsection{SHAP}
Por otro lado, SHAP \cite{xai_shap} es una técnica explicativa que se basa en la teoría de juegos cooperativos \cite{juegos_cooperativos}. SHAP asigna a cada variable un valor de importancia conocido como Shapley value, que representa su contribución a la predicción. A diferencia de LIME, que se centra en explicaciones locales para predicciones individuales, SHAP permite tanto el análisis local de predicciones concretas como la visión global de la importancia de las variables en todo el conjunto de datos, ofreciendo propiedades deseables como consistencia, precisión local y aditividad. En el contexto de este trabajo, la combinación de SHAP y LIME permitirá identificar qué variables macroeconómicas, como los tipos de interés, la inflación, el diferencial de la curva de tipos o el precio del petróleo, son más influyentes en las predicciones del rendimiento de los fondos de inversión, proporcionando así una comprensión completa y robusta del comportamiento del modelo.

\section{Trabajos anteriores}
Por último, para concluir con el marco teórico, se verán trabajos relacionados con este TFG, con el objetivo de contextualizar la investigación actual dentro del panorama más amplio de la aplicación de Inteligencia Artificial en finanzas. Y es que, en los últimos años, el uso de técnicas de Machine Learning y Deep Learning en el ámbito financiero ha crecido muy significativamente, debido a la gran cantidad de datos y el fácil acceso a los mismos, a su vez, los avances computacionales, también han impulsado este crecimiento. Tal y como se menciona en la revisión sistemática de El Alami y colaboradores \cite{elalami2025}, los modelos más utilizados en el ámbito financiero son Random Forest, XGBoost, Support Vector Machines, y redes neuronales como LSTM, Bidirectional LSTM y CNN, del mismo modo, se hace uso de modelos híbridos que combinan métodos estadísticos con inteligencia artificial.

En cuanto al rendimiento de estos modelos, tal y como señalan Nabipour Mojtaba y colaboradores \cite{nabipour2020}, las arquitecturas basadas en LSTM superan consistentemente a los métodos tradicionales como ARIMA o GARCH en la predicción de series financieras, esto es gracias a su capacidad para capturar dependencias temporales largas. No obstante, los propios autores advierten que las LSTM pueden ser propensas al sobreajuste y requieren un ajuste cuidadoso de sus hiperparámetros, algo que se tendrá en cuenta en la fase de modelado de este trabajo.

En lo que respecta al uso de variables macroeconómicas, tal y como explica Abdou Housein en su artículo “The impact of oil and global markets on Saudi stock market predictability: A machine learning approach” \cite{abdou2024}, el hecho de introducir variables como los precios del petróleo, los tipos de interés y los índices de volatilidad mejora significativamente la predictibilidad de los mercados bursátiles cuando se incorporan dentro de marcos de Machine Learning. Este artículo, da aún más fuerza a la decisión adoptada en este TFG de incluir un importante número de variables macroeconómicas como predictores del rendimiento de los fondos de inversión o activos financieros.

En el aspecto de la interpretabilidad de los modelos, es un tema que ha ganado mucha popularidad en los últimos años, esto queda reflejado en el artículo de Mohsin y Nasim “Explaining the Unexplainable: A Systematic Review of Explainable AI in Finance.” \cite{mohsin2025}, donde se explica que la falta de interpretabilidad de los modelos complejos limita su adopción en entornos regulados como el financiero, donde la transparencia es clave. Para solucionar este problema, Ye y Chen \cite{ye2025}, proponen técnicas como SHAP y LIME, que permiten explicar las predicciones de modelos de caja negra, demostrando que es posible realizar modelos complejos manteniendo la transparencia necesaria.

Por último, tal y como se señala en la revisión sistemática sobre forecasting financiero \cite{vancsura2025}, actualmente, muchos estudios evalúan sus modelos solo con métricas estadísticas sin integrar las predicciones en estrategias de inversión reales, a su vez pocos, modelos son validados en diferentes ciclos económicos como pueden ser crisis financieras o períodos de alta inflación. Estas limitaciones refuerzan el enfoque adoptado en este TFG, que combina la predicción multivariante-multistep con un análisis exhaustivo de interpretabilidad mediante SHAP y LIME, y una evaluación en diferentes contextos económicos identificados mediante clustering.
